% 本文件由 LaTeX 排版 Agent 根据有效 translation.json 自动生成。
% author=Tertullian; work_id=0301; title=辩护

\chapter{辩护}

% structure: p0001 source=h1 target=omitted reason=page-title-already-rendered-by-parent

% structure: p0002 source=h2 target=section reason=relative-heading-rank; base=h2

\section{第一章}

罗马帝国的统治者们，如果你们坐在高高的审判席上施行公义，在众目睽睽之下，且几乎占据国家最高地位，却不能公开调查并当众审查针对基督徒的指控之真相；如果唯独在此事上，你们害怕或羞于以符合正义的审慎态度行使权力进行公开调查；如果最近在私下判决中对我们的子民施加的极端严厉，阻碍了我们得以在你们面前为自己辩护，那么你们必定无法阻止真理通过无声著作的秘密路径到达你们的耳中。真理就她的处境而言，无需向你们上诉，因为那并不令她惊异。她知道自己在世上不过是寄居者，且身在陌生人中自然发现仇敌；不仅如此，她的起源、居所、希望、赏报、尊荣都在天上。与此同时，她只焦急地向世间统治者祈求一事——不被未经审理即定罪。让律法在其领域内至高无上，给予她一次听审，有何害处？甚至，岂不更显其绝对权威：即使在她提出辩护后仍定她的罪？但如果未听审就对她宣判，那么除了不公正行为的憎恶之外，你们还将招致理所当然的怀疑，即你们这样做带有某种不公正的念头，因为不愿听取你们可能无法听完并定罪的话。我们向你们提出这一点，作为我们主张你们对基督徒之名的仇恨是不公正的第一理由。而似乎为这种不公正开脱的理由（我指无知）恰恰同时加重并证明了不公正。因为还有什么比憎恨你一无所知的事物更不公平，即使它该被憎恨？只有当那憎恨被\Emph{已知}为应得时，它才配得。但若没有知识，它的正义如何得以证明？因为正义要从对主题的了解来证明，而非仅仅从存在厌恶这一事实。当人们仅仅因为完全不知道所厌恶事物的性质而任由厌恶时，为什么那恰恰可能正是他们不应厌恶的事物？所以我们主张，他们既因无知而恨我们，又因继续无知而不义地恨我们，二者互为因果。证明他们无知（同时谴责并开脱其不公正）的是：那些曾经因对基督教一无所知而恨它的人，一旦认识它，便立刻放下敌意。从仇敌变为门徒。仅仅通过了解，他们开始憎恨自己从前的样子，并信奉从前所恨的；其人数之多正如加诸我们头上的指控。呼声是：国家充满了基督徒——他们在田野、在堡垒、在岛屿上；他们哀叹，好像遭了灾祸，说两性、每个年龄和身份，甚至高位者，都皈依了基督信德；然而尽管如此，他们的心智仍未觉醒，去思考其中未被注意到的善。他们不允许任何更真实的怀疑进入脑海；他们不愿更仔细地检验。唯独在此，人类的好奇心入睡了。他们乐于无知，尽管对他人而言知识已是福乐。\Person{阿那卡西斯}责备鲁莽者妄图批判有教养者；他对这种完全无知之人评判有知之人，更会怎样痛斥！因为他们已经厌恶，便不想知道更多。因此他们预先判断他们无知的东西，认为假若他们认识它，它便不能再是厌恶的对象；因为如果调查发现没有可厌之处，当然应停止不公的厌恶；而如果它的坏性质清楚显现，那么厌恶非但不会减少，反而在正义本身的权威下获得了坚持这种厌恶的更强理由。但有人说，一件事物好，不是因为多数人皈依它；因为有多少人天性倾向恶！有多少人误入歧途！这无疑是的。然而，彻底邪恶的事物，甚至那些被它裹挟的人也不敢辩护它是好的。本性给一切邪恶披上恐惧或羞愧的面纱。例如，你们发现罪犯急于隐藏自己，避免公开露面，被捕时惶恐不安，受审时否认罪行；即使受酷刑，也并非轻易或总是承认；当定罪毫无疑义时，他们为自己的行为悲伤。在内心反省中，他们承认自己受邪恶倾向驱使，却归咎于命运或星辰。他们不愿承认这是他们自己的事，因为他们承认这是邪恶的。但在基督徒的情况中有什么类似？他感到的唯一羞愧或遗憾，是没有更早成为基督徒。若被指出，他引以为荣；若被控告，他不辩护；被审问，他自愿承认；被判罪，他感谢。这是什么样邪恶的事，竟然缺乏邪恶所有通常的特性——恐惧、羞愧、托词、悔恨、哀叹？什么！那是一桩罪行，而罪犯竟为之欢喜？被人指控是他热切期望，因被指控而受罚是他的幸福？你们不能称它为疯狂，你们这些被证明对此事一无所知的人。\par

% structure: p0004 source=h2 target=section reason=relative-heading-rank; base=h2

\section{第二章}

若再假定我们确是最恶之人，为何你们对待我们与对待其他罪犯如此不同？即对待同类罪犯，理应同罪同罚。当他人被控时，允许他们用自己口舌或雇佣辩护人证明清白；他们有充分机会答辩和辩论；事实上，法律禁止未经辩护和听审就定任何人的罪。唯独基督徒被禁止为自身辩解，为真理辩护，帮助法官作出公正判决；所关心的只是满足公众仇恨所要求的事——承认此名，而非审查指控。而在你们通常的司法调查中，若有人承认杀人、渎神、乱伦或叛国（这正是我们被控的罪名），你们不会满足于立即判决——直到彻底审查供认的细节——行为的真实性质、次数、地点、方式、何时、共谋者、实际参与者，才会采取这一步。在我们案件中，毫无类似做法，尽管关于我们的谎言本应同样被筛查，以便查明我们每人尝过多少被杀孩童的肉，每人暗中犯过多少乱伦，有何厨子、何狗见证我们的行为。哦，若能查出某个吞食百名婴孩的基督徒，那将是统治者多大的荣耀！但相反，我们发现甚至对我们的调查也被禁止。因为\Person{小普林尼}在担任行省总督时，曾处死一些基督徒，并迫使一些坚忍之人动摇，但仍为人数众多而烦恼，最后向在位皇帝\Person{图拉真}请示如何处理其余人，向他的主上说明：除了顽固不愿献祭外，他在宗教仪式中只发现清晨聚会唱诗赞美基督和天主，并通过共同承诺忠于宗教来封印其生活方式，禁止杀人、奸淫、盗窃及其他罪行。对此，\Person{图拉真}回信说：基督徒绝不应被搜捕；但若被带到法庭前，则应受惩罚。啊，可悲的谕令——在案情必要下，自相矛盾！它禁止搜捕他们，视之为无辜；又命令惩罚他们，视之为有罪。它既仁慈又残忍；既放过又惩罚。审判啊，你为何对自己玩弄回避的把戏？你若定罪，为何不先调查？你若不调查，为何不宣告无罪？行省各处设有军事哨所追捕强盗。对付叛徒和公敌，人人皆为士兵；甚至搜索其同谋和从犯。唯独基督徒不得被搜捕，尽管他可被带到法官面前受控；仿佛搜索还有其他目的！因此，你们判处此人，无人曾想搜捕他，他却出现在你面前，而即使现在他也不应受罚，我想，并非因其罪过，而是因为虽禁止搜捕，却被发现。而且，你们处理我们时并不遵循针对罪犯的常规司法程序；因为对于其他否认者，你们施以酷刑逼供；唯独对基督徒，你们施以酷刑逼其否认；而若我们犯有任何罪行，我们定会否认，你们则会用酷刑强迫我们招供。的确，你们不应仅因我们承认此名就认定罪行无需调查——\Emph{你们}日常虽深知何为谋杀，却仍从已认罪的杀人犯口中逼取犯罪经过的全部详情。因此，你们行事更加荒谬：当我们因承认基督之名而被你们认定有罪时，你们却用酷刑迫使我们背弃承认，以便在否认此名的同时，也否认你们曾依据同一承认所假定我们有过的罪行。我想，虽然你们认为我们是最坏的人，却不愿我们灭亡。因为无疑，你们惯于命令杀人犯否认，若亵渎神明者坚持认罪则将其送上刑架！是这样吗？但若你们如此对待我们，不视我们为罪犯，你们就宣告我们无罪，因为作为无罪之人，你们急于让我们不要坚持一个你们知道会使我们遭受你们手中必然而非公正的定罪的承认。那人喊道：\Quote{我是基督徒。}他告诉你他是什么；你希望从他口中听到他不是什么。你身居权威之位，为逼出真相，却尽力从我们这里得到谎言。他说：\Quote{我是}你所问我是不是的。为何你折磨我犯罪？我承认，你却施以酷刑。若我否认，你又会怎样？当然你们并不轻易相信其他否认者。当我们否认时，你们立即相信。让你们的这种荒谬行为使你们怀疑此案中有某种隐藏的力量，在它的影响下，你们违背常规、违背公义的本质，甚至违背法律本身行事。因为，除非我大错特错，法律命令搜捕罪犯，而非隐藏他们。法律规定认罪者应被定罪，而非无罪释放。元老院法令、长官命令都明确如此。你们所服务的权力是文明的，而非暴虐的统治。的确，暴君曾将酷刑用作惩罚；在你们，则减轻为仅为审讯手段。坚持你们的法律，在必要情况下直至取得供认；若酷刑因供认而提前，则无需实施；应作出判决，将罪犯交付应得刑罚，而非释放。因此，无人急于让有罪者无罪开释；那样做不正当，故无人被迫否认。然而，你们认为基督徒是各种罪行之人，是诸神、皇帝、法律、良善道德、整个自然的敌人；却强迫他否认，以便你宣告他无罪——若他不否认，你们无法做到。你们行事反复无常。你们希望他否认罪行，以便即使违背他意愿，你们也使他成为清白且与过去一切罪责无关的人！你们这奇怪的荒谬从何而来？你们为何不反思：自愿的承认远比被迫的否认更可信；或考虑一下，当被迫否认时，一个人的否认是否真诚，无罪开释后，他是否立即在审判结束后嘲笑你们的敌意，依然如故是基督徒？既然你们在各方面对待我们与其他罪犯不同，一心只想夺走我们的名号（确实，若我们做了基督徒从不做的事，这名号就不再属于我们），那就完全清楚，此案中没有任何罪行，只是一个名字，一个常与真理作对的体系以敌意追逼它，主要目的是确保人们不愿确切知道他们确知完全不知之事。因此，他们也相信关于我们他们无法证实的事，不愿深入调查，以免他们宁可信以为真的指控全被证明毫无根据，以免那敌视敌对势力的名字——其罪行被假定而非证实——仅凭自己承认就被定罪。所以我们若承认就受酷刑，若坚持就受惩罚，若否认就被释放，因为所有争议只关乎一个名字。最后，为何你们从名簿中读出某人是个基督徒？为何不也说他是个杀人犯？若基督徒是杀人犯，岂不也犯乱伦或你们相信我们的其他恶行？唯独在我们身上，你们羞于或不愿提及我们罪行的名称——若被称为\Quote{基督徒}不意味着任何罪行，那么这名字本身就极为可憎，当它本身就被当作罪行时。\par

% structure: p0006 source=h2 target=section reason=relative-heading-rank; base=h2

\section{第三章}

我们该如何看待这件事：多数人如此盲目地撞向对基督徒名号的仇恨，以致当他们为某人作有利证词时，总要混入对他所持名号的辱骂？有人说：\Quote{一个好男人，就是\Person{盖乌斯·塞伊乌斯}，只因为他是一个基督徒。}另一个人则说：\Quote{我很惊讶像\Person{卢奇乌斯}这样明智的人竟会突然成了基督徒。}没有人认为有必要考虑：盖乌斯是否不好、卢奇乌斯是否不智，正因他是基督徒；或者一个人是基督徒，正因他明智良善。他们赞美所知的，辱骂所不知的，并以无知启发其知识；尽管公平说来，你应当由已知判断未知，而非由未知判断已知。另一些人，对于他们认识为放荡、卑鄙、邪恶的人，在他们取了基督徒名号之前，却从他们所赞美的事上加以烙印。在仇恨的盲目中，他们抨击自己赞同的判断！\Quote{那女子曾是何等一个人！多么淫荡！多么轻佻！那青年曾是何等一个人！多么放荡！多么好色！——他们成了基督徒！}于是，可恨的名号被赋予品格的改变。有些人甚至以舒适换取那仇恨，甘愿承受伤害，只要家中能远离他们深恶痛绝的对象。妻子现今贞洁了，丈夫不再嫉妒，却将她逐出家门；儿子现今顺从了，素来忍耐的父亲却剥夺其继承权；仆人现今忠诚了，一度温和的主人却命令他离开眼前；任何人因这被憎恶的名号而改过迁善，都成了大罪。善良的价值不如对基督徒的仇恨。好，若对这名号有如此厌恶，你们能归咎于名号什么罪过？你们能控告纯粹的称号什么，除非词中某音听来野蛮、不祥、粗鄙或淫秽？但基督徒一词，就字义而言，源于傅油。是的，甚至当你们错误念作\Quote{Chrestianus}时（因为你们连所恨的名号都不确切知道），它也来自甘饴与良善。因此，你们在无罪者身上，甚至恨一个无罪的名称。但憎恶这教派的特殊理由是，它以其创立者之名命名。一个宗教教派从其宗师得名于追随者，这有什么新奇呢？哲学家们不也是从\Emph{其}学派的创始人得名——柏拉图学派、伊壁鸠鲁学派、毕达哥拉斯学派吗？斯多葛学派和学院派不也是从他们聚集和驻留的地方得名吗？医生不也是从\Person{埃拉西斯特拉图斯}得名，文法学家从\Person{阿里斯塔尔库斯}得名，甚至厨子从\Person{阿皮基乌斯}得名吗？然而，带着从最初创立者连同其所创立的一切传下来的名号，并不冒犯任何人。毫无疑问，如果证明这教派是邪恶的，那么其创立者同样邪恶，这便将证明名号是恶的，应受我们厌恶，考虑到教派及其创始人的品性。因此，在憎恶这名号之前，你们理应在创始人中考察教派，或在教派中考察创始人。但如今，未经任何审察和认识二者，单是名号就成了控诉的缘由，单是名号就受攻击，单是一个声音就给教派及其创始人两者带来定罪，而你们对两者都一无所知，只因为他们有这样一个名称，并非因他们被判有罪。\par

% structure: p0008 source=h2 target=section reason=relative-heading-rank; base=h2

\section{第四章}

既然我以上所说等于序言，为要揭露公众对我们仇恨的不义，我现在就要立足在我们无瑕的辩护上；我不仅要驳斥加于我们的指控，还要将这些指控反过来加在指控者身上，这样所有人都能知道基督徒远非他们深知盛行于他们自己中间的罪行，同时使他们因指控我们而脸红——我并非说最坏的人控告最好的人，而是按他们现在所愿，控告那些不过是与他们同罪的人。我们将回应各种暗中被控的罪行，而这些罪行我们却发现他们在光天化日之下犯下；我们因这些罪行而被视为邪恶、无知、该受惩罚、该受嘲笑。但是，当我们的真理在各方面战胜你们时，法律权威作为最后手段被搬出来反对它，以至于要么说法律的决定是绝对结论性的，要么，尽管不情愿，服从的必要性被置于真理之上。因此，我首先要在这法律问题上，与你们作为法律选定的维护者进行辩论。首先，当你们严厉地在判决中宣布说：\Quote{你们存在是不合法的}，并以毫不犹豫的严厉命令执行这一点时，你们表现出了纯粹暴政的暴力和不义的统治，如果你们仅仅因为希望它不合法就否认某事合法，而不是因为它本当不合法。但如果你们认为它不合法是因为它\Emph{本当}不合法，那么无疑那些有害的事不应得到法律许可；而基于此，事实上，已经确定凡是有益的就是合法的。那么，如果我发现了你们法律所禁止的事是好的，作为已有此先见的人，这法律不已经失去了阻止我的效力吗？尽管那同一件事，如果是邪恶的，它就会公正地禁止我。如果你们的法律错了，我想，它是源于人的，并非从天降下。人会在立法上犯错，或在废弃法律时醒悟，这有什么奇怪吗？难道斯巴达人不是修改了\Person{来古格士}本人的法律，从而对其作者造成如此痛苦，以致他把自己关起来，绝食至死吗？你们自己不是每天都在努力照亮古代的黑暗，用皇帝诏令和敕令的新斧头砍削你们法律那古老粗糙的整片森林吗？难道最坚定的统治者\Person{塞维鲁}不是昨天才废除了可笑的帕皮亚法吗？该法强迫人们在犹利亚法允许缔结婚姻之前就生育子女，尽管它们有年代的权威。古时也有法律，规定败诉的一方可能被债权人肢解；然而，后来因普遍同意，这种残忍从法令中删除，死刑改为耻辱标记。通过采用没收债务人财产的办法，旨在让他的脸红得出血，而不是让他流血。有多少法律仍隐藏在视线之外需要改革！因为法律受称赞既非由于年代久远，也非由于立法者的尊贵，而仅仅因为它们公正；因此，当它们的不义被认识到时，即使它们定罪，也理当被谴责。我们为何说它们不义？不，如果它们惩罚纯粹的名称，我们大可以称它们为不理性的。但如果它们惩罚行为，为什么在我们的事例中，它们仅凭名称就惩罚行为，而在其他事例中，它们必须从行为本身而非从名称来证明？我是行乱伦的人（他们这样说），为何他们不调查？我是杀婴者，为何他们不用酷刑从我这里获取实情？我犯了对神明、对凯撒的罪，为何我这能为自己辩白的人，不被允许为自己辩护？没有法律禁止审查它所禁止的罪行，因为法官若不确信罪行已发生，就绝不会施行正义的惩罚；公民若不知道所受惩罚之事的性质，也就不会真正服从法律。法律公正还不够，法官相信其公正也不够；那些被期待服从的人也应同样信服。不，法律若不在乎被检验和认可，就处于强烈嫌疑之下：未经证实就统治人，这绝对是邪恶的法律。\par

% structure: p0010 source=h2 target=section reason=relative-heading-rank; base=h2

\section{第五章}

要谈一谈我们现在所提及的这类法律的起源，曾有一条古老的法令：未经元老院批准，皇帝不得封立任何神祇。\Person{马尔库斯·埃米利乌斯}在其神\Person{阿尔布尔努斯}的事上就有此经验。而这一点也有利于我们的案件，因为你们是按人的判断来赋予神性。除非诸神使人满意，否则它们就不能被神化：神必须取悦于人。因此，\Person{提庇留}——在他执政期间基督徒的名号进入了世界——从\Place{巴勒斯坦}得到了那些明确显示基督神性真实性的消息，便亲自将此事提交元老院，并提出了支持基督的裁决。元老院因未亲自批准，否决了他的提议。\Person{恺撒}坚持自己的意见，威胁要愤怒地对待所有控告基督徒的人。查考你们的历史；你们会发现在那里\Person{尼禄}是第一个用帝国之剑攻击基督徒教派的人，当时尤其在罗马取得进展。但我们却因受如此恶徒的敌意而使我们的定罪得以圣化而自豪。因为知道尼禄的人都能理解，除非一件事有非凡的卓越之处，否则不会招致尼禄的定罪。还有\Person{多米提安}，一个在残酷上与尼禄同类型的人，也尝试过迫害；但由于他尚存一丝人性，很快就终止了他所开始的事，甚至召回了那些他曾流放的人。这样的人一直是我们迫害者——不义、不敬、卑鄙，连你们自己也无从称赞他们，通常你们会恢复那些受他们判决之人的自由。然而，从那时直到现在，在如此众多拥有天上人间智慧的君主中，请指出一个迫害基督徒之名的。绝非如此，相反，我们向你们引介一位保护者，正如你们通过查考最庄重的皇帝\Person{马尔库斯·奥勒留}的书信所看到的，他在信中作证说，日耳曼旱灾因那些恰巧在他麾下作战的基督徒的祷告所获得的雨水而解除了。他虽然没有通过公法解除基督徒的法律限制，却以另一种方式公开将其搁置一旁，甚至对他们的控告者增加了更严厉的定罪判决。这些是什么法律？只有不敬、不义、卑贱、血腥、麻木、疯狂的人才执行它们来反对我们。\Person{图拉真}在一定程度使其作废，禁止搜寻基督徒；而\Person{哈德良}——尽管他热衷于探求一切奇异新奇之事——\Person{韦斯巴芗}——尽管他是犹太人的征服者——\Person{庇护}——\Person{维鲁斯}——都从未执行过这些法律。更合理的判断应是：恶人理当被善良的君王——他们的天敌——所根除，而不是被与自身气质相似的人所根除。\par

% structure: p0012 source=h2 target=section reason=relative-heading-rank; base=h2

\section{第六章}

我现在要请这些最虔诚的祖制法律的保护者和辩护者告诉我，论到他们自己对祖传制度的忠诚、尊敬和遵从，他们是否从未偏离——是否从未离开古道——是否从未摒弃那作为道德生活准则中最有用和最必要的一切。那些抑制奢侈和炫耀生活方式的律法如今何在？它们禁止一餐花费超过一百\OriginalTerm{阿司}{asses}，一次只可上一只禽鸟，且不得是肥的；它们曾以一项视为重大的理由将一位贵族元老逐出元老院，因为他拥有了十磅白银，因而图谋不轨；它们曾在戏院刚兴起败坏民风时便迅速将其取缔；它们不容许官职尊严或高贵出身的徽章被轻率或不受惩罚地篡用。因为我看到如今的\OriginalTerm{百钱宴}{Centenarian suppers}名称所由来的，已不是一百阿司，而是\OriginalTerm{塞斯特斯}{sestertia}的开销；银矿被制成餐具（若这仅适用于元老，而不适用于被释奴甚至只配鞭打的贱民，那还算小事）。我也看到，一座戏院不够，且戏院都不加屋顶：无疑是为了使不贞的享乐在冬季也不致沉闷，拉栖代蒙人发明了供戏剧用的羊毛斗篷。我现在看到主妇与妓女的服饰已无区别。关于妇女，你们祖宗的律法曾如此鼓励端庄与节操，现在也已废弃；那时妇女除了手指上的——即丈夫以婚戒神圣许给自己——之外，身上尚无黄金；妇女戒酒如此严格，以致一位主妇因打开酒窖隔间而被亲友饿死——而在\Person{罗慕路斯}时代，仅仅因尝酒，\Person{梅塞尼乌斯}就杀死了妻子，且未因此受惩罚。为此，妇女有亲吻亲属的习俗，以便通过气息被发觉。那令人如此向往的婚姻生活的幸福何在？它曾是我们早期风俗的特色，其结果是在约六百年间我们中间没有一次离婚。如今，妇女全身上下挂满黄金；酗酒如此普遍，以致亲吻从不出于自愿；至于离婚，她们渴望之，如同婚姻的自然结果。你们祖宗的智慧所颁布的关于神祇本身的律法，你们这些最忠诚的子孙也已废除。执政官奉元老院之权，将\Person{巴克斯}神及其秘仪不仅逐出城市，而且逐出整个意大利。执政官\Person{皮索}和\Person{加宾纽斯}（当然不是基督徒）禁止\Person{塞拉皮斯}、\Person{伊西斯}、\Person{阿波克拉蒂斯}连同他们的犬首神进入\Place{卡皮托山}——以此将他们逐出众神集会——推倒他们的祭坛，将他们驱逐出境，唯恐他们卑劣纵欲宗教的恶行蔓延。你们却恢复了这些，并授予他们最高荣誉。你们对祖宗教的尊敬何在？在服饰、饮食、生活方式、见解，甚至最后在言语上，你们都背弃了祖先。你们总是赞美古风，却每天在生活方式上花样翻新。从你们未能持守应守的一切，你们清楚表明：你们一面抛弃祖先的良好作风，一面保留并守护着不该保留的东西。然而，你们似乎仍如此忠实地捍卫、并据以主要控告基督徒的祖传传统——我指崇拜神祇的热忱，即古代所犯的主要错误——虽然你们重建了如今已成为罗马神祇的塞拉皮斯的祭坛，并向如今已成为意大利神的巴克斯献上狂欢祭，我将在适当的地方指明你们是如何蔑视、忽略并推翻古人的权威，完全弃之不顾的。同时，我将继续回应那关于秘密罪行的恶名指控，为我探讨公开之事扫清道路。\par

% structure: p0014 source=h2 target=section reason=relative-heading-rank; base=h2

\section{第七章}

我们这些邪恶的怪物被人指控举行一种神圣的仪式，在其中杀死一个小孩子并吃掉；在宴席之后，我们行乱伦，那些狗——我们的拉皮条者，确实，打翻灯火，使我们在黑暗的无耻中满足不敬的欲望。这就是常被加于我们的指控，而你们却不费心去查明我们长久被控告之事的真相。那么，如果你们相信，就将这事带到光天化日之下；若从未查究，就不要相信它。基于你们的反复无常，我们有权向你们断言，你们不敢\OriginalTerm{查明}{expiscate}的事，必无真实。在基督徒的案件上，你们加给刽子手一种与\OriginalTerm{查明}{expiscation}完全相反的职责：他并非要他们承认自己所行，而是要他们否认自己的身份。我们如前所述，将我们宗教的起源追溯到\Person{提庇留}皇帝在位之时。真理与对真理的憎恨一同进入世界。真理一出现，就被视为仇敌。它有同样多的敌人：犹太人——不出所料，出于竞争之心；士兵们——出于勒索钱财的欲望；甚至我们的家仆——出于他们的本性。我们每日被敌人围困，每日被出卖；我们常在聚会和集会上被突袭。谁曾按照那流行的说法，遇到一个啼哭的婴儿？谁曾将独眼巨人和塞壬的血盆大口原封不动地保留给法官？谁曾在他们妻子身上发现任何不洁的痕迹？哪里有人发现了如此暴行却隐瞒起来，或者在将罪犯拖到法官面前时，因受贿而保持沉默？如果我们总是保守秘密，我们的行为何时为世人所知？不，又能被谁知晓呢？当然不是由犯罪者自己；因为对于秘仪，忠诚的沉默永远是义务。\OriginalTerm{萨莫色雷斯秘仪}{Samothracian}和\OriginalTerm{厄琉西斯秘仪}{Eleusinian}从不泄露——更何况对于这些一旦揭露就立即招致人间惩罚、而天主的忿怒被留存到将来的事，沉默岂不更应被遵守？那么，如果基督徒不是自己散布自己的罪行，那当然一定是外人。而他们从哪里得知呢？因为在宗教入门仪式中，普遍习俗是阻止非信徒靠近，并提防见证人，除非那些如此邪恶的人比他们的邻人更少畏惧？每个人都知道谣言是什么。这是你们自己的谚语之一：\Quote{在一切恶事中，没有比谣言传得更快的。}为什么谣言是如此的恶事？是因为它迅速？是因为它传递信息？还是因为它极端虚假？——即便它给我们带来一些真理，也无不夹杂虚假，要么减少，要么增加，要么改变了简单的事实。不仅如此，它的存在法则就是只在撒谎时持续，只在无证据时生存；因为一旦有了证据，它就不再存在；而它完成了仅仅传播谣言的使命之后，便交出事实，从此被视为事实，被称为事实。于是没有人会说，例如，\Quote{他们说这事发生在罗马}，或\Quote{有谣言说他获得了一个行省}，而是说\Quote{他得到了一个行省}和\Quote{这事发生在罗马。}谣言，即不确定的名称，在事情确定时便无立足之地。除了愚人，谁会相信它？因为智者从不相信可疑之事。每个人都知道，无论它多么热心地被传播，无论它基于多么强烈的断言，它总有某时从某个源头产生。然后它必须爬入传播的舌头和耳朵；一个小小的初始瑕疵如此污染了整个故事，以至于无人能确定最初发出它的嘴唇，是否播下了虚假的种子——像常发生的那样，出于对立精神，或出于多疑的判断，或出于一种根深蒂固的、甚至在某些人身上是天生的、对撒谎的喜爱。好在时间使一切水落石出，正如你们的谚语和格言所证明的，这是天性的安排，它如此规定，使任何事都无法长久隐藏，即使谣言并未散布它。那么，长久以来只有名声知晓基督徒的罪行，这正该如此。这就是你们用来反对我们的见证——这个见证从未能证明它某时散布的指控，而最终仅凭持久性变成了世上固定的看法；因此我自信地诉诸永远真实的天性本身，反驳那些毫无根据地认为这些事应当被相信的人。\par

% structure: p0016 source=h2 target=section reason=relative-heading-rank; base=h2

\section{\Emph{第八章}}

现在，你看，我们将这些滔天罪行的报酬摆在你面前。它们应许永生。请暂且以此为你的信念。那么我问你，既然这样相信，你是否认为以你将要拥有的良心来获得它是值得的？来，将你的刀刺入那婴孩，他无仇敌，无罪指控，是众人的孩子；或者若那是他人的事，你只需置身于一个尚未真正活着便死去的人身旁，等候那刚赋予的灵魂离去，接住那鲜嫩的血液，用它浸透你的面包，自由地享用。当你斜倚在桌前时，留意你的母亲和妹妹所坐的位置；记清楚，以便当狗造成的黑暗降临在你身上时，你不会犯错误，因为你将犯下罪行——除非你做出乱伦之事。经过如此入门与封印，你便获得永生。请你告诉我，永生值得吗？若不值得，则这些事不可信。即使你相信，我也否认你的意愿；即使你有意愿，我也否认可能性。那么，若你不能，别人为何能？若别人能，你为何不能？我想我们是不同的本性——我们是\OriginalTerm{狗头人}{Cynopæ}还是\OriginalTerm{独脚人}{Sciapodes}？你本人和基督徒一样是人：若你做不到，就不该相信别人能做到，因为基督徒也和你一样是人。但无知之人，确实受欺骗和愚弄。他们完全不知道有这类事情被归咎于基督徒，否则他们一定会亲自调查了解。相反，我想，那些希望入教的人的习惯是，首先去找主持者，请他说明需要做哪些准备。那么，在这种情况下，他无疑会说：\Quote{你必须有一个尚在幼年、不知死为何物、能在你的刀下微笑的孩子；还要有面包，用以收集喷涌的血；此外，还需要烛台、灯和狗——以及引诱它们扑灭灯光的小块食物；最重要的是，你必须带上你的母亲和妹妹。}但若母亲和妹妹不愿意呢？或者根本没有母亲和妹妹呢？若基督徒没有基督徒亲属呢？我想，一个没有兄弟或儿子的人，不会被认为是基督的真正门徒。那么，如果他们事先不知道这些事，却正在等待着他们呢？至少事后他们会知道；他们忍受着并宽恕。据说，他们害怕若泄露秘密会付出代价：不，他们反而会因此有充分的理由得到保护；人们或许认为，他们甚至宁愿自杀，也不愿背负如此可怕的知识而活着。就算他们有这种恐惧，但他们为何还坚持呢？因为很明显，如果你事先知道这些，你就绝不会愿意继续那些你原本就不会做的事。\par

% structure: p0018 source=h2 target=section reason=relative-heading-rank; base=h2

\section{第九章}

为了更彻底地驳斥这些指控，我将指出，在你们中间，部分公开、部分秘密地盛行着一些做法，这些做法或许使你们也相信了关于我们的类似事情。在提比略担任总督期间，直到最近，在非洲还公开向农神献祭儿童，他将祭司们悬挂在遮蔽神庙的圣树上示众——如此多的十字架，正义所要求的惩罚在这些十字架上追上了他们的罪行，我们国家的士兵至今仍可作证，他们正是为那位总督做那项工作的。甚至现在，这种神圣的罪行仍在秘密进行。你们看，不仅仅是基督徒鄙视你们；因为你们所做的一切，没有任何罪行被彻底而持久地根除，你们的任何神祇也没有改变他的作风。当农神不饶恕自己的儿女时，他也不可能饶恕别人的儿女；事实上，正是这些儿女的父母自己习惯于献祭他们，欣然回应加于他们的召唤，并在那种场合下使小孩子们高兴，以免他们流着泪死去。同时，杀人与弑亲之间有天壤之别。在高卢，更年长的年龄被献祭给墨丘利。我把陶里克人的传说交给他们自己的剧院。为什么，即使在虔诚的埃涅阿斯后裔的那座最宗教性的城市里，也有某个朱庇特，在他们的游戏中人们用人的血来洗他。你说，那是斗兽士的血。难道因为它是一个邪恶之人的血，就变得不那么是人的血了吗？还是更低贱？无论如何，它是在谋杀中流出的。哦，朱庇特，你自己就是一个基督徒，而且实在是你残忍父亲的独生子！但是关于杀婴，无论它是为神圣目的还是仅仅出于个人冲动而犯下的——尽管我们说过，弑亲与杀人大有不同——我将转向一般民众。想想看，那些拥挤着、张着嘴渴望基督徒之血的人——甚至你们的许多统治者，他们以对你们的正义和对我们的严厉措施而著称——有多少我可以指控他们凭良心犯了杀害自己后代的罪？至于杀人的种类，用淹死、或暴露于寒冷、饥饿和狗来杀死当然是更残忍的方式。更成熟的年龄总是更喜欢刀剑之死。在我们这里，杀人一旦被绝对禁止，我们甚至不能毁掉子宫中的胎儿，尽管此时人尚从身体其他部位获取血液以维持生命。阻止出生只是更快的杀人；无论你是夺走已出生的生命，还是毁灭即将出生的生命，都没有区别。那个将要成为人的，已经是人了；你已经在种子中拥有了果实。至于血餐和诸如此类的悲惨菜肴，请读一读——我不确定在哪里记载（我想是在希罗多德中）——从手臂上取血，并由双方尝过，如何成为某些民族之间的条约纽带。我不确定在喀提林时代尝过的是什么。他们还说到，在一些斯基泰部落中，死者被他们的朋友吃掉。但我离题太远了。就在今天，在我们中间，献给贝娄娜的血，从刺穿的大腿上抽出然后被分食，封印了进入那位女神仪式的入门。那些在角斗士表演中，为了治疗癫痫，贪婪地渴饮竞技场上被杀死的罪犯伤口中流出的鲜血，然后匆忙跑开的人——他们属于谁？那些在战斗地点以野兽的肉为餐的人——谁对熊和鹿有强烈的食欲？那只熊在搏斗中沾满了它撕裂之人的血；那只鹿在角斗士的血泊中打滚。熊的内脏里满是尚未消化的人的内脏，被大肆需求。你们有人打着嗝，吃的是人喂养的肉？如果你们享用这样的食物，你们的宴席与我们基督徒被指控的宴席有什么不同？而那些以野蛮的欲望夺取人体的人，他们吞噬活人，难道做得更少吗？他们因为只舔食将要变成血的东西，就较少沾染人血的污秽吗？他们显然吃的不是婴儿，而是成年人。为你们卑劣的行径在基督徒面前脸红吧，他们即使在简单自然的食物餐中也没有动物的血；他们禁食被勒死和自然死亡的东西，唯一的原因是为了不沾染污秽，如同不沾染藏在内脏里的血一样。用单一的例子来结束此事，你们用血肠诱惑基督徒，正是因为你们完全知道，你们试图让他们违反的这件事，他们认为是非法的。认为那些你们确信连尝牛血都感到恐惧的人，却渴望人血，这是多么不合理啊；除非，也许你们尝试过，发现人血更甜美！事实上，这里有一个你们可以用来发现基督徒的测试，就像火盆和香炉一样。他们应该通过他们对人血的胃口来证明，就像通过他们拒绝献祭来证明一样；否则，他们应该通过拒绝尝血来确认为与基督教无关，就像通过献祭一样；而人血不会缺乏，因为在审判和定罪囚犯时会大量提供。那么，还有谁比那些受过朱庇特本人教诲的人更倾向于乱伦的罪行呢？克特西亚斯告诉我们，波斯人与他们的母亲有非法性关系。马其顿人也在这点上受到怀疑；因为在第一次听到俄狄浦斯的悲剧时，他们嘲笑了乱伦者的悲伤，喊道，\par

\SourceRecord{Translated by S. Thelwall.；Ante-Nicene Fathers；Vol. 3.；Edited by Alexander Roberts, James Donaldson, and A. Cleveland Coxe.；Buffalo, NY: Christian Literature Publishing Co.,；1885.；<http://www.newadvent.org/fathers/0301.htm>.}
